\fchapter{Ej 2 Pág 65}

\fsection{\boldit{\textcolor{waveBlue2}{Investiga}}{Como ves,
hay muchos tipos de modelos as a Service. Aunque IaaS, PaaS, SaaS y XaaS son los más conocidos, hay algunos 
otros muy interesantes en el mercado que pueden ofrecer soluciones más específicas a nuestro negocio. 
Busca en internet información sobre CaaS, SECaaS, NaaS, DBaaS y DRaaS, y comenta las características de 
cada uno de ellos con el resto de la clase.}}

\begin{solution}
	Aqui tenemos resumidamente las características de cada uno de los modelos as a Service:
	\begin{itemize}
		\item \boldit{CaaS (Containers as a Service):} Ofrece una plataforma en la nube para desplegar y gestionar aplicaciones en contenedores (como Docker), eliminando la necesidad de administrar la infraestructura subyacente.

		\item \boldit{SECaaS (Security as a Service):} Consiste en la externalización de servicios de ciberseguridad a un proveedor en la nube, ofreciendo protección (como firewalls, antivirus o detección de intrusiones) bajo un modelo de suscripción.

		\item \boldit{NaaS (Network as a Service):} Permite consumir servicios de red (como conectividad, banda ancha o firewalls) a través de la nube, sin necesidad de adquirir y mantener hardware de red físico.

		\item \boldit{DBaaS (Database as a Service):} Es un servicio en la nube que proporciona bases de datos completamente gestionadas, incluyendo aprovisionamiento, copias de seguridad, parches y escalado automático.

		\item \boldit{DRaaS (Disaster Recovery as a Service):} Ofrece una solución completa de recuperación ante desastres en la nube, replicando sistemas y datos para permitir una restauración rápida en caso de una interrupción o fallo.
	\end{itemize}

\end{solution}
